\chapter{Obtener el código}
\label{cap:preparacion}

Todos los modos parten de una copia del repositorio en el servidor, siempre en la misma carpeta:

\begin{xltabular}{\linewidth}{@{}l L@{}}
	\toprule
	\textbf{Servidor} & \textbf{Carpeta del código} \\
	\midrule
	\endhead
	Linux & \cod{/opt/digaipe/app} \\
	Windows & \texttt{C:\bs digaipe\bs app} \\
	\bottomrule
\end{xltabular}

En Docker, \texttt{docker compose} se ejecuta desde la subcarpeta \texttt{deploy} de esa copia;
en servidor físico, la aplicación se ejecuta directamente desde ella.

\section{Instalar Git}

\so{Linux}
\begin{codigo}
sudo apt update
sudo apt install -y git
\end{codigo}

\so{Windows} Descargue e instale Git for Windows desde \texttt{https://git-scm.com/download/win}
con las opciones por defecto, o, si el servidor tiene \texttt{winget}:
\begin{codigo}
winget install --id Git.Git -e --source winget
\end{codigo}
Cierre y vuelva a abrir PowerShell para que se actualice la variable \texttt{PATH}.

\section{Clonar el repositorio}

\so{Linux}
\begin{codigo}
sudo mkdir -p /opt/digaipe
sudo git clone https://github.com/luishok/digaipe.git /opt/digaipe/app
cd /opt/digaipe/app
sudo git checkout <RAMA_O_ETIQUETA_A_INSTALAR>
sudo git log -1 --format="%h %cd"
\end{codigo}

La copia pertenece a \texttt{root}; por eso todos los comandos \texttt{git} sobre ella llevan
\texttt{sudo} (sin él, Git se niega con \emph{detected dubious ownership}).

\so{Windows}
\begin{codigo}
New-Item -ItemType Directory -Force C:\digaipe
git clone -c core.autocrlf=false https://github.com/luishok/digaipe.git C:\digaipe\app
Set-Location C:\digaipe\app
git checkout <RAMA_O_ETIQUETA_A_INSTALAR>
git log -1 --format="%h %cd"
\end{codigo}

Si el repositorio es privado, Git pedirá usuario y un \emph{token} de acceso de GitHub (no la
contraseña de la cuenta). Use un token de solo lectura limitado a este repositorio.

Anote en la hoja de datos el commit que muestra el último comando.

\begin{aviso}[Finales de línea]
	Los scripts \texttt{.sh} deben tener finales de línea LF; con CRLF fallan con
	\texttt{\$'\textbackslash r': command not found}. En Windows, Git convierte por defecto a CRLF:
	por eso se clona con \texttt{-c core.autocrlf=false}. Compruébelo en ambos sistemas:
	\begin{codigo}
git ls-files --eol deploy/backup/ scripts/docker/
	\end{codigo}
	(en Linux, con \texttt{sudo} delante)
	La segunda columna debe decir \texttt{w/lf} en todos los archivos.
\end{aviso}

\begin{nota}
	La copia del servidor es solo para instalar: no edite archivos dentro de ella. La
	configuración (contraseñas) va en archivos que Git ignora (\texttt{deploy/.env}) o fuera de la
	carpeta (servidor físico), y los archivos cargados están en \texttt{var/uploads}, que Git
	también ignora. Así, actualizar es un \texttt{git pull} sin conflictos
	(capítulo~\ref{cap:actualizacion}).
\end{nota}
